\section{术语消化：本期关键词索引}

\begin{longtable}{p{0.24\textwidth}p{0.34\textwidth}p{0.33\textwidth}}
\toprule
术语 & 一句话解释 & 在本期中的作用 \\
\midrule
Physical AI & 物理世界中的 AI，能感知、决策并执行动作 & 小鹏从自动驾驶升级到 AI 战略的核心框架。 \\
Robotaxi & 自动驾驶出租车服务 & 刘先明希望在广州跑好的阶段性目标。 \\
Software 1.0 & 规则、优化和手写逻辑主导的软件栈 & 早期自动驾驶路线。 \\
Software 2.0 & 神经网络和数据驱动的软件栈 & 端到端自动驾驶的基础概念。 \\
VLM & Vision-Language Model，视觉语言模型 & 理解视觉和语言，但不一定直接输出 action。 \\
VLA & Vision-Language-Action，视觉语言动作模型 & 小鹏 Physical AI 技术栈的重要方向。 \\
Language Bottleneck & 把连续传感器信息压成语言 token 的中间瓶颈 & 刘先明主张拆掉的关键中间层。 \\
Self-supervised Learning & 自监督学习，用数据自身构造训练信号 & 支撑大规模数据利用，减少人工标注依赖。 \\
Scaling & 通过更大模型、更多数据、更多算力提升能力 & 刘先明认为 Physical AI 的关键 bet。 \\
云端模型工厂 & 云端训练大模型，再压缩部署到车端 & 连接 scaling 和车端量产的工程结构。 \\
蒸馏 & 用大模型训练小模型，让小模型继承能力 & 云端到车端部署的重要手段。 \\
量化/剪枝 & 降低模型精度或裁剪参数以适配硬件 & 车端部署和成本控制需要。 \\
Infra & 数据、训练、分析和部署基础设施 & Cruise 和小鹏路线共同强调的底座。 \\
Corner Case & 长尾边界场景 & 自动驾驶数据闭环要持续挖掘的对象。 \\
\bottomrule
\end{longtable}

\subsection{本章小结}

本期术语围绕一个问题展开：如何把自动驾驶从规则工程升级为可 scaling 的 Physical AI 系统。关键词不是孤立名词，而是训练、压缩、部署和量产反馈的一条链。

